% Bloque independiente. Sustituir los campos por evidencia de su proyecto.
\bloque{\personaCuatro}
\section{Llamadas y fallos de red}
\begin{frame}{Llamadas y fallos de red}\relax
{\large\bfseries Una llamada bloqueante espera la respuesta. Una API no bloqueante permite continuar al hilo cliente. El plazo de espera delimita cuánto esperar.\par}
\medskip
{\small Fuente de apoyo \cite{kazanavicius2023,rozek2026}.}
\guia{\campo{Explicar un concepto con sus palabras y relacionarlo con la demo.}}
\note{Completar el guion en notas/secciones/04-integrante-4.tex.}
\end{frame}

\begin{frame}{Ejemplo sencillo}\relax
\textbf{Operación}\quad \campo{Una llamada RPC/gRPC/RMI con respuesta y un fallo controlado}\par
\medskip
\textbf{Entrada}\quad \campo{Dato de prueba}\par
\textbf{Salida}\quad \campo{Resultado realmente obtenido}\par
\guia{Declarar la tecnología elegida y mostrar qué recibe el cliente cuando el servicio falla.}
\end{frame}

\begin{frame}{\proyectoCuatro · arquitectura propuesta}\relax
{\small\textbf{Actualmente}\quad \campo{Problema, usuarios, componentes y cifra con fuente o medición.}}
\par\medskip
\centering
\begin{tikzpicture}
\node[nodo] (a) {Cliente\\{\campo{Componente actual}}};
\node[nodo,right=3cm of a] (b) {Servicio propuesto\\{\campo{Operación remota}}};
\draw[enlace] (a)--node[above,font=\small]{\campo{Protocolo y datos}}(b);
\end{tikzpicture}
\par\medskip
\raggedright
{\small\textbf{Propuesta}\quad \campo{Qué se separaría y por qué la comunicación elegida aporta valor.}}
\guia{Redibujar a partir de su bosquejo a mano. Marcar límites de proceso, red, datos y elementos que ya existen.}
\end{frame}

\begin{frame}{Decisiones al distribuir el proyecto}\relax
\small
\begin{tabularx}{\linewidth}{@{}lX@{}}
\toprule
Comunicación & \campo{Operación, contrato y llamada bloqueante o no bloqueante. Justificar.}\\
Transparencias & \campo{Una ofrecida y una limitada, con ejemplo para el usuario.}\\
CAP & \campo{Topología supuesta, enlace que se corta y conducta de lectura/escritura.}\\
PACELC & \campo{Con la red sana, qué consistencia y latencia buscaríamos y cómo medirlas.}\\
\bottomrule
\end{tabularx}
\guia{Vincular cada decisión con el flujo del diagrama. Precisar réplicas y garantías antes de asignar una etiqueta CAP.}
\end{frame}

\begin{frame}{Fallo hipotético y lección para el proyecto}\relax
\small
\textbf{Escenario supuesto}\quad \campo{Durante una operación, el servicio o enlace falla.}\par
\textbf{Causa y recuperación}\quad \campo{Cadena causal, detección y respuesta propuesta.}\par
\textbf{Comparación}\quad \campo{Sistema real comparable, fuente y compromiso técnico.}\par
\textbf{Lección}\quad \campo{Qué implementaríamos primero y con qué prueba lo validaríamos.}\par
\guia{Separar el resultado esperado de una prueba ya realizada. Si prueban el fallo, registrar pasos, salida y fecha como simulación.}
\end{frame}
